% GENERATED FILE. Edit canonical lessons, then run npm run generate:books.
% canonical-source-hash: fnv1a64:c990c6aa16431c92
% canonical-lessons: MR-C03-kasa, MR-C03-tumhi-kase-aahat, MR-C03-mi, MR-C03-mi-bara-aahe, MR-C03-kaahi-harkat-nahi, MR-C03-practice

\chapter{How Are You?}
\label{ch:responding}
\hlchaptermodality{10}

\begin{chapteropening}
\textbf{By the end of this chapter:} \emph{I can ask how someone is in Marathi, answer naturally, and return the question.}

Complete MR-C03-practice, the canonical closing checkpoint, and demonstrate the chapter promise: ask how someone is in Marathi, answer naturally, and return the question.
\end{chapteropening}

\section[kasā / kaśī / kasaṁ]{\mr{कसा} / \mr{कशी} / \mr{कसं} (kasā / kaśī / kasaṁ) — \textquotedblleft{}how\textquotedblright{}}

\label{lesson:MR-C03-kasa}



\begin{quote}
The everyday \textquotedblleft{}how are you?\textquotedblright{} needs one more \emph{k-} question word — \textbf{how}. And, being an adjective-like word, it too carries gender.
\end{quote}

\subsection*{The letters in this word}
\emph{(Skim if you read Devanagari.)} \textbf{\mr{क}} (\emph{ka}) + \textbf{\mr{सा}} (\emph{sā}) $\to$ \textbf{\mr{कसा}} \emph{kasā}; \textbf{\mr{क}} + \textbf{\mr{शी}} (\emph{śī}) $\to$ \textbf{\mr{कशी}} \emph{kaśī}; \textbf{\mr{क}} + \textbf{\mr{सं}} (\emph{saṁ}) $\to$ \textbf{\mr{कसं}} \emph{kasaṁ}.

\begin{cousinweb}
\textbf{\mr{कसा}} (\emph{kasā}, \textquotedblleft{}how\textquotedblright{}) belongs to the same \textbf{k-} question family as \emph{kāy} (what), \emph{koṇ} (who), \emph{kuṭhe} (where) — all from PIE \textbf{*k\textsuperscript{w}o-}, the root behind English \textbf{how} (yes, \emph{how} is a \emph{wh-} word that lost its \emph{w}-spelling) and Latin \emph{quo-modo} (\textquotedblleft{}in what way,\textquotedblright{} which became Spanish \emph{cómo}).
\end{cousinweb}

\begin{grammarlens}[title={Grammar Lens: \textquotedblleft{}how\textquotedblright{} agrees with who you ask}]
Because it describes the person's state, \emph{how} takes their gender:

\begin{itemize}
\raggedright
  \item \textbf{\mr{कसा}} (\emph{kasā}) — asking a \textbf{man}
  \item \textbf{\mr{कशी}} (\emph{kaśī}) — asking a \textbf{woman}
  \item \textbf{\mr{कसं}} (\emph{kasaṁ}) — asking about a neuter thing
\end{itemize}

So \textquotedblleft{}how are you?\textquotedblright{} changes shape with whom you address — you will see this in the next lesson.
\end{grammarlens}

\subsection*{Your turn}
Say these aloud:

\begin{itemize}
\raggedright
  \item \textquotedblleft{}kasā / kaśī / kasaṁ\textquotedblright{} — how, in three genders
  \item two cousins from the same root (English \emph{how}, Spanish \emph{cómo})
\end{itemize}

\begin{tcolorbox}[breakable,colback=teal!4,colframe=teal!35!black,title={Before you move on}]
Which \textquotedblleft{}how\textquotedblright{} do you use asking a woman? (\emph{kaśī} — feminine.)
\end{tcolorbox}
\section[tumhī kase āhāt?]{\mr{तुम्ही} \mr{कसे} \mr{आहात}? (tumhī kase āhāt?) — \textquotedblleft{}how are you?\textquotedblright{}}

\label{lesson:MR-C03-tumhi-kase-aahat}



\begin{quote}
The full courteous \textquotedblleft{}how are you?\textquotedblright{} — assembled from words you already hold.
\end{quote}

\subsection*{You'll want to know — the phrase, assembled}
\textbf{\mr{तुम्ही}} (\emph{tumhī}, respectful \textquotedblleft{}you\textquotedblright{}) + \textbf{\mr{कसे}} (\emph{kase}, \textquotedblleft{}how,\textquotedblright{} plural-agreeing) + \textbf{\mr{आहात}} (\emph{āhāt}, \textquotedblleft{}are,\textquotedblright{} respectful) $\to$ \textbf{\mr{तुम्ही} \mr{कसे} \mr{आहात}?} — \textquotedblleft{}you how are?\textquotedblright{}, verb last.

Two agreements are at work:

\begin{itemize}
\raggedright
  \item Because \textbf{\mr{तुम्ही}} is grammatically plural, \textquotedblleft{}how\textquotedblright{} appears as \textbf{\mr{कसे}} (\emph{kase}, the plural form of \emph{kasā}), and the copula as \textbf{\mr{आहात}} (\emph{āhāt}).
  \item To a single male friend: \textbf{\mr{तू} \mr{कसा} \mr{आहेस}?} (\emph{tū kasā āhes?}); to a female friend: \textbf{\mr{तू} \mr{कशी} \mr{आहेस}?} (\emph{tū kaśī āhes?}).
\end{itemize}

\begin{culture}
The respectful \emph{tumhī … āhāt} is the safe default with anyone new. The gendered \emph{kasā/kaśī} only appears in the familiar \emph{tū} form — one more place Marathi's three genders surface in ordinary speech.
\end{culture}

\subsection*{Your turn}
Say these aloud:

\begin{itemize}
\raggedright
  \item \textquotedblleft{}tumhī kase āhāt?\textquotedblright{} — how are you? (respectful)
  \item it to a female friend (\emph{tū kaśī āhes?})
\end{itemize}

\begin{tcolorbox}[breakable,colback=teal!4,colframe=teal!35!black,title={Before you move on}]
Why \emph{kase} and \emph{āhāt} (not \emph{kasā}, \emph{āhes}) with \emph{tumhī}? (\emph{tumhī} is grammatically plural, so \textquotedblleft{}how\textquotedblright{} and \textquotedblleft{}are\textquotedblright{} take plural forms.)
\end{tcolorbox}
\section[mī]{\mr{मी} (mī) — \textquotedblleft{}I\textquotedblright{}}

\label{lesson:MR-C03-mi}



\begin{quote}
To answer \textquotedblleft{}how are you?\textquotedblright{}, you need to say \textbf{I}. In Marathi that is a single bright syllable: \textbf{\mr{मी}} (\emph{mī}).
\end{quote}

\subsection*{The letters in this word}
\emph{(Skim if you read Devanagari.)} \textbf{\mr{मी}} = \textbf{\mr{म}} (\emph{ma}) + the long-\emph{ī} stroke \textbf{\mr{ी}} $\to$ \emph{mī}.

\begin{cousinweb}
\textbf{\mr{मी}} (\emph{mī}, \textquotedblleft{}I\textquotedblright{}) descends from Sanskrit \textbf{\mr{अहम्}} (\emph{aham}, \textquotedblleft{}I\textquotedblright{}), which wore down over centuries to the short \emph{mī}. That \emph{aham} comes from PIE \textbf{*h\textsubscript{1}eǵ(H)om} — the ancestor of Latin \textbf{ego} and English \textbf{I}. And notice its family tie to last chapter's \textbf{\mr{माझं}} (\emph{mājhaṁ}, \textquotedblleft{}my\textquotedblright{}): the same first-person \textbf{m-} thread runs through \textquotedblleft{}I\textquotedblright{} and \textquotedblleft{}my,\textquotedblright{} just as English pairs \textbf{I} with \textbf{me/my}.
\end{cousinweb}

\begin{grammarlens}[title={Grammar Lens: \textquotedblleft{}I am\textquotedblright{} is \emph{mī āhe}}]
Pair \textbf{\mr{मी}} with the copula \textbf{\mr{आहे}}: \textbf{\mr{मी} \mr{आहे}} (\emph{mī āhe}, \textquotedblleft{}I am\textquotedblright{}). To answer that you are well, reach back to Chapter 4's \textbf{\mr{बरं}} (\emph{baraṁ}, \textquotedblleft{}fine\textquotedblright{}) in its gendered form: \textbf{\mr{मी} \mr{बरा} \mr{आहे}} (\emph{mī barā āhe}, a man) / \textbf{\mr{मी} \mr{बरी} \mr{आहे}} (\emph{mī barī āhe}, a woman) — \textquotedblleft{}I am fine.\textquotedblright{}
\end{grammarlens}

\subsection*{Your turn}
Say these aloud:

\begin{itemize}
\raggedright
  \item \textquotedblleft{}mī\textquotedblright{} — I; then \textquotedblleft{}mī āhe\textquotedblright{} — I am
  \item \textquotedblleft{}I am fine\textquotedblright{} for yourself (\emph{mī barā/barī āhe})
\end{itemize}

\begin{tcolorbox}[breakable,colback=teal!4,colframe=teal!35!black,title={Before you move on}]
What links \emph{mī} (\textquotedblleft{}I\textquotedblright{}) and \emph{mājhaṁ} (\textquotedblleft{}my\textquotedblright{})? (The shared first-person \textbf{m-} root — as English links \emph{I} with \emph{me/my}.)
\end{tcolorbox}
\section[mī barā āhe]{\mr{मी} \mr{बरा} \mr{आहे} (mī barā āhe) — \textquotedblleft{}I'm well, thank you\textquotedblright{}}

\label{lesson:MR-C03-mi-bara-aahe}



\begin{quote}
The answer to \textquotedblleft{}how are you?\textquotedblright{} — and it reuses a word you already banked in Chapter 4.
\end{quote}

\subsection*{You'll want to know — the phrase, assembled}
\textbf{\mr{मी}} (\emph{mī}, I) + \textbf{\mr{बरा}} (\emph{barā}, \textquotedblleft{}well/fine\textquotedblright{}) + \textbf{\mr{आहे}} (\emph{āhe}, am) $\to$ \textbf{\mr{मी} \mr{बरा} \mr{आहे}} — \textquotedblleft{}I am well.\textquotedblright{} Add thanks: \textbf{\mr{मी} \mr{बरा} \mr{आहे}, \mr{धन्यवाद}} (\emph{mī barā āhe, dhanyavād}).

Back in Chapter 4 you met \textbf{\mr{बरं}} (\emph{baraṁ}) standing alone as \textquotedblleft{}okay/fine\textquotedblright{} — its \textbf{neuter} form. When it describes \emph{you}, it takes your gender:

\begin{itemize}
\raggedright
  \item \textbf{\mr{बरा}} (\emph{barā}) — a man is well
  \item \textbf{\mr{बरी}} (\emph{barī}) — a woman is well
  \item \textbf{\mr{बरं}} (\emph{baraṁ}) — the bare \textquotedblleft{}okay\textquotedblright{} (neuter)
\end{itemize}

\begin{grammarlens}[title={Grammar Lens: the copula changes with the subject}]
You now have three forms of \emph{āhe} in play: \textbf{\mr{मी} \mr{आहे}} (\emph{mī āhe}, I am), \textbf{\mr{तू} \mr{आहेस}} (\emph{tū āhes}, you are, familiar), \textbf{\mr{तुम्ही} \mr{आहात}} (\emph{tumhī āhāt}, you are, respectful). One verb, reshaped for each person — Marathi marks \emph{who} on the verb itself.
\end{grammarlens}

\subsection*{Your turn}
Say these aloud:

\begin{itemize}
\raggedright
  \item \textquotedblleft{}mī barā āhe, dhanyavād\textquotedblright{} (man) / \textquotedblleft{}mī barī āhe, dhanyavād\textquotedblright{} (woman)
  \item the three copula forms — \emph{mī āhe, tū āhes, tumhī āhāt}
\end{itemize}

\begin{tcolorbox}[breakable,colback=teal!4,colframe=teal!35!black,title={Before you move on}]
Why does the bare \emph{baraṁ} you banked earlier become \emph{barā} here? (It now describes a person, so it takes that person's gender.)
\end{tcolorbox}
\section[kāhī harkat nāhī]{\mr{काही} \mr{हरकत} \mr{नाही} (kāhī harkat nāhī) — \textquotedblleft{}no problem / you're welcome\textquotedblright{}}

\label{lesson:MR-C03-kaahi-harkat-nahi}



\begin{quote}
When someone thanks you, Marathi's easy reply is not \textquotedblleft{}you're welcome\textquotedblright{} but \textquotedblleft{}there's \textbf{no objection} at all.\textquotedblright{}
\end{quote}

\subsection*{You'll want to know — the phrase, assembled}
\textbf{\mr{काही}} (\emph{kāhī}, \textquotedblleft{}any\textquotedblright{}) + \textbf{\mr{हरकत}} (\emph{harkat}, \textquotedblleft{}objection, trouble\textquotedblright{}) + \textbf{\mr{नाही}} (\emph{nāhī}, \textquotedblleft{}not/none\textquotedblright{}) $\to$ \textquotedblleft{}any objection there-is-not\textquotedblright{} $\approx$ \textbf{\textquotedblleft{}no problem at all.\textquotedblright{}}

\begin{cousinweb}
Two threads meet here:

\begin{itemize}
\raggedright
  \item \textbf{\mr{नाही}} (\emph{nāhī}, \textquotedblleft{}not\textquotedblright{}) is Chapter 4's \textquotedblleft{}no,\textquotedblright{} on the ancient negative \textbf{na-} — PIE \textbf{*ne}, the root of English \textbf{no, not, none} and Latin \emph{nōn}.
  \item \textbf{\mr{हरकत}} (\emph{harkat}, \textquotedblleft{}objection, hindrance\textquotedblright{}) is a \textbf{Perso-Arabic} loan (from Arabic \emph{ḥarakah}), one of the many words Marathi absorbed over centuries of contact with Persian-speaking rulers and traders — above all under the \textbf{Deccan sultanates} — the same trade-and-court layer that colours all the languages of the Deccan.
\end{itemize}

So the reply blends an Indo-European negative with an Arabic noun — a small map of Marathi's history in three words.
\end{cousinweb}

\begin{culture}
Where English \textquotedblleft{}thank you / you're welcome\textquotedblright{} trades politeness tokens, Marathi waves the thanks away: \emph{don't mention it, there was no trouble.} The same gesture, a different picture.
\end{culture}

\subsection*{Your turn}
Say these aloud:

\begin{itemize}
\raggedright
  \item \textquotedblleft{}dhanyavād\textquotedblright{} $\to$ \textquotedblleft{}kāhī harkat nāhī\textquotedblright{} — thanks $\to$ no problem
  \item the two histories inside it (\emph{nāhī} $\leftarrow$ PIE \emph{ne; }harkat\emph{ $\leftarrow$ Arabic)}
\end{itemize}

\begin{tcolorbox}[breakable,colback=teal!4,colframe=teal!35!black,title={Before you move on}]
What two language families meet in \emph{kāhī harkat nāhī}? (Indo-European \emph{nāhī} $\leftarrow$ \emph{*ne; Arabic }harkat\emph{.)}
\end{tcolorbox}
\section[Practice]{Chapter 10 recap — how are you}

\label{lesson:MR-C03-practice}



\begin{quote}
The whole everyday exchange, start to finish.
\end{quote}

\subsection*{The exchange}
\noindent
\begin{tabularx}{\linewidth}{@{}>{\raggedright\arraybackslash}X>{\raggedright\arraybackslash}X>{\raggedright\arraybackslash}X@{}}
\toprule
\textbf{Marathi} & \textbf{Say it} & \textbf{English} \\
\midrule
\mr{तुम्ही} \mr{कसे} \mr{आहात}? & \emph{tumhī kase āhāt?} & How are you? \\
\mr{मी} \mr{बरा} \mr{आहे}, \mr{धन्यवाद}. & \emph{mī barā āhe, dhanyavād.} & I'm well, thank you. \\
\mr{काही} \mr{हरकत} \mr{नाही}. & \emph{kāhī harkat nāhī.} & No problem / you're welcome. \\
\bottomrule
\end{tabularx}

\begin{grammarlens}[title={What you've built — atoms banked this chapter}]
\begin{itemize}
\raggedright
  \item \textbf{\mr{कसा} / \mr{कशी} / \mr{कसं}} (\emph{kasā/kaśī/kasaṁ}) — how; the \emph{k-} question family, gender-agreeing.
  \item \textbf{\mr{मी}} (\emph{mī}) — I; Sanskrit \emph{aham} $\to$ Latin \textbf{ego}, English \textbf{I}.
  \item \textbf{\mr{बरा} / \mr{बरी} / \mr{बरं}} (\emph{barā/barī/baraṁ}) — well/fine, now gendered.
  \item \textbf{\mr{काही} \mr{हरकत} \mr{नाही}} (\emph{kāhī harkat nāhī}) — no problem; \emph{nāhī} $\leftarrow$ PIE \emph{*ne}, \emph{harkat} $\leftarrow$ Arabic.
  \item Copula forms: \textbf{\mr{मी} \mr{आहे} / \mr{तू} \mr{आहेस} / \mr{तुम्ही} \mr{आहात}} (\emph{mī āhe / tū āhes / tumhī āhāt}) — one verb, reshaped per person.
\end{itemize}
\end{grammarlens}

\subsection*{Your turn}
Say these aloud:

\begin{itemize}
\raggedright
  \item the full three-line exchange
  \item \textquotedblleft{}how are you?\textquotedblright{} and \textquotedblleft{}I'm well\textquotedblright{} for your own gender
\end{itemize}

\begin{tcolorbox}[breakable,colback=teal!4,colframe=teal!35!black,title={Before you move on}]
Give the respectful \textquotedblleft{}how are you?\textquotedblright{} and a full answer. (\emph{tumhī kase āhāt?} — \emph{mī barā/barī āhe, dhanyavād.}) Next chapter: the farewells — \emph{punhā bheṭū.}
\end{tcolorbox}
